\section*{Chapter \ref{ch:b1:first-law} --- The First Law of Thermodynamics}

\begin{solution}{exo:b1:first-law:1}
$Q = C_{P,m}\Delta T = 29.1 \times 100 = \qty{2.91}{kJ}$; $W = -nR\Delta T =
\qty{-0.83}{kJ}$; $\Delta U = Q + W = \qty{2.08}{kJ} = C_{V,m}\Delta T$.
\end{solution}

\begin{solution}{exo:b1:first-law:2}
$W = -nRT\ln(V_2/V_1) = -2494\ln(0.25) = \qty{+3.46}{kJ}$; $\Delta U = 0$ so
$Q = \qty{-3.46}{kJ}$ (given to the bath).
\end{solution}

\begin{solution}{exo:b1:first-law:3}
$T_2 = T_1 \times 10^{\gamma - 1} = 300 \times 10^{0.4} = \qty{754}{K}$; $P_2 =
10^{1.4} = \qty{25}{bar}$.
\end{solution}

\begin{solution}{exo:b1:first-law:4}
Argon: $C_{V,m} = \tfrac32 R = 12.5$, $C_{P,m} = \tfrac52 R = 20.8$ J/(mol K),
$\gamma = 1.67$; nitrogen: $20.8$, $29.1$, $\gamma = 1.40$. For water the
expansion on heating is tiny, so the work $P\,\dd V$ at constant pressure
is negligible against the \omterm{def:b1:first-law:heat}{heat}: $c_P - c_V \approx 0$.
\end{solution}

\begin{solution}{exo:b1:first-law:5}
$T_f = (200 \times 80 + 300 \times 20)/500 = \qty{44}{\degreeCelsius}$. With the
cup: $(200 \times 4.18 \times 80 + (300 \times 4.18 + 80) \times 20)/(500 \times
4.18 + 80) = \qty{43}{\degreeCelsius}$.
\end{solution}

\begin{solution}{exo:b1:first-law:6}
Ice: warm to $0$: $0.050 \times 2100 \times 10 = \qty{1.05}{kJ}$; melt:
\qty{16.7}{kJ}; total \qty{17.8}{kJ}. Water can give \qty{25.1}{kJ} down to
$0$: all melts; then $0.250 \times 4180\,T_f = 25.1 - 17.8$ kJ: $T_f =
\qty{7}{\degreeCelsius}$. With \qty{150}{g}: needs \qty{53.3}{kJ} $>$ \qty{25.1}{kJ}:
only $(25.1 - 3.15)/334 = \qty{66}{g}$ of ice melt, $T_f = \qty{0}{\degreeCelsius}$,
\qty{84}{g} of ice remain.
\end{solution}

\begin{solution}{exo:b1:first-law:7}
$W = -P_1(V_2 - V_1) + 0 - P_2(V_1 - V_2) + 0 = (P_2 - P_1)(V_2 - V_1) > 0$:
the cycle is traversed counterclockwise (expansion at low pressure,
compression at high), so the gas \emph{receives} net work, equal to the
enclosed area; clockwise it would deliver it --- an engine.
\end{solution}

\begin{solution}{exo:b1:first-law:8}
$T_2 = 300 \times 7^{0.286} = \qty{523}{K}$. The hot compressed air \omterm{def:b1:first-law:heat}{heats} the
barrel by conduction at each stroke. In the tire the air cools to
ambient at fixed volume: $P \propto T$, the pressure falls by the ratio
of temperatures --- top it up.
\end{solution}

\begin{solution}{exo:b1:first-law:9}
$W = 0$ (no external pressure to push against), $Q = 0$: $\Delta U = 0$,
hence $\Delta T = 0$ for a \omterm{def:b1:kinetic-theory:model}{perfect gas}. Irreversible: the gas never flows
back by itself. A real gas, whose molecules attract, cools slightly
(part of the \omterm{thm:b1:work-and-energy:ke}{kinetic energy} becomes \omterm{def:b1:work-and-energy:conservative}{potential energy} as they separate).
\end{solution}

\begin{solution}{exo:b1:first-law:10}
$Q = L_v = \qty{2.26}{MJ}$; $W = -P(V_g - V_l) = -1.013\times10^5 \times 1.67 =
\qty{-0.17}{MJ}$; $\Delta U = \qty{2.09}{MJ}$: $93\%$ of the \omterm{def:b1:first-law:heat}{heat} goes into
separating the molecules (\omterm{def:b1:kinetic-theory:U}{internal energy}), $7\%$ into pushing back the
atmosphere.
\end{solution}

\begin{solution}{exo:b1:first-law:11}
$W = -P_2(V_2 - V_1)$, $\Delta U = nC_{V,m}(T_2 - T_1)$, $V_2 = nRT_2/P_2$, $V_1
= nRT_1/P_1$: $\tfrac52(T_2 - 300) = -T_2 + 300 \times 5$, $3.5T_2 = 2250$,
$T_2 = \qty{643}{K}$, $V_2 = \qty{10.7}{L}$. Reversible: $T_2 = 300 \times
5^{0.286} = \qty{475}{K}$, $V_2 = \qty{7.9}{L}$. The sudden compression does
more work on the gas (the full \qty{5}{bar} acts from the start) and ends
hotter and less compressed.
\end{solution}

\begin{solution}{exo:b1:first-law:12}
$\rho = PM/RT = \qty{1.20}{kg/m^3}$: $\sqrt{\gamma P/\rho} = \sqrt{1.4 \times
1.013\times10^5/1.20} = \qty{343}{m/s}$; $\sqrt{P/\rho} = \qty{290}{m/s}$, $15\%$
low. With $P/\rho = RT/M$: $c = \sqrt{\gamma RT/M}$; against $u = \sqrt{3RT/M}$
= \qty{502}{m/s}: $c/u = \sqrt{\gamma/3} = 0.68$.
\end{solution}

\begin{solution}{pb:b1:first-law:1}
\textbf{1.} $n = P_1V_1/RT_1 = 10^5 \times 2.12\times10^{-4}/2494 = \qty{8.5e-3}{mol}$.

\textbf{2.} $PV^\gamma$ constant: $V_2 = V_1(P_1/P_2)^{1/\gamma} = 212 \times
7^{-0.714} = \qty{53}{cm^3}$.

\textbf{3.} $T_2 = T_1(P_2/P_1)^{(\gamma - 1)/\gamma} = 300 \times 7^{0.286} =
\qty{523}{K}$.

\textbf{4.} $W = nC_{V,m}(T_2 - T_1) = 8.5\times10^{-3} \times 20.8 \times 223 =
\qty{39}{J}$.

\textbf{5.} Air to add: $\Delta n = \Delta P\,V/RT = 6\times10^5 \times 2\times10^{-3}/
2494 = \qty{0.48}{mol}$: $0.48/8.5\times10^{-3} = 57$ strokes, about
\qty{2.2}{kJ} --- a minute of honest effort.

\textbf{6.} At constant volume $P \propto T$: the pressure falls as the air
cools (by up to $300/523$ for the last stroke's air); the cyclist adds a
few strokes after a pause.

\textbf{7.} $W_{\text{iso}} = nRT\ln(P_2/P_1) = 8.5\times10^{-3} \times 2494 \times
1.95 = \qty{41}{J} > \qty{39}{J}$: the isothermal path compresses the gas
further (to \qty{30}{cm^3} instead of \qty{53}{cm^3}) to reach \qty{7}{bar},
because it stays cool --- more volume swept, more work.

\textbf{8.} $n = 10^5 \times 5\times10^{-4}/2494 = \qty{0.020}{mol}$; $T_2 =
300 \times 20^{0.4} = \qty{994}{K}$.

\textbf{9.} $P_2 = 20^{1.4} = \qty{66}{bar}$.

\textbf{10.} \qty{994}{K} is well above the fuel's self-ignition point: the
injected fuel lights on contact. A petrol engine compresses its
air--fuel mixture only to about \qty{750}{K} so that it does not ignite
before the spark; too high a ratio gives ``knock''.

\textbf{11.} $W = nC_{V,m}(T_2 - T_1) = 0.020 \times 20.8 \times 694 = \qty{289}{J}$.

\textbf{12.} $25 \times 289 = \qty{7.2}{kW}$ per cylinder, stored in the hot
gas and given back in the expansion.

\textbf{13.} Irreversibility (the gas does not follow the quasi-static
path; extra work is dissipated) raises $T_2$; \omterm{def:b1:first-law:heat}{heat} lost to the walls
lowers it. In a real engine the second usually wins: $T_2$ is a little
below the ideal value.

\textbf{14.} $Q = 20\times10^{-6} \times 43\times10^6 = \qty{860}{J}$.

\textbf{15.} Isobaric: $Q = nC_{P,m}(T_3 - T_2)$: $T_3 - T_2 = 860/(0.020 \times
29.1) = \qty{1480}{K}$, $T_3 = \qty{2470}{K}$; $V_3 = V_2T_3/T_2 = 25 \times
2.48 = \qty{62}{cm^3}$.

\textbf{16.} $W = -P_2(V_3 - V_2) = -66\times10^5 \times 37\times10^{-6} =
\qty{-244}{J}$: work delivered by the gas.

\textbf{17.} $T_4 = T_3(V_3/V_1)^{\gamma - 1} = 2470 \times (0.124)^{0.4} =
\qty{1070}{K}$; $W = nC_{V,m}(T_4 - T_3) = 0.020 \times 20.8 \times (-1400) =
\qty{-582}{J}$.

\textbf{18.} Net work delivered $= 244 + 582 - 289 = \qty{537}{J}$; ratio
$537/860 = 62\%$ (the ideal Diesel efficiency at these settings; real
engines reach about $40\%$).

\textbf{19.} $1 - 300/2470 = 88\%$: the ideal Diesel cycle stays well
below the Carnot bound, because its \omterm{def:b1:first-law:heat}{heat} is received over a range of
temperatures rather than at the highest one.

\textbf{20.} Over a cycle $\Delta U = 0$: $Q_{\text{out}} = -(860 - 537) =
\qty{-323}{J}$, i.e.\ \qty{323}{J} leave with the exhaust; check:
$nC_{V,m}(T_4 - T_1) = 0.020 \times 20.8 \times 770 = \qty{320}{J}$.

\textbf{21.} The oscillation is fast (a second) compared with \omterm{def:b1:first-law:heat}{heat}
conduction through \qty{10}{L} of air: adiabatic. Linearizing
$PV^\gamma$: $\dd P/P = -\gamma\,\dd V/V = -\gamma Sx/V$.

\textbf{22.} Force on the ball $S\,\dd P = -\gamma PS^2x/V$: $m\ddot x =
-(\gamma PS^2/V)x$, harmonic, $\omega_0^2 = \gamma PS^2/mV$.

\textbf{23.} $\omega_0^2 = 1.4 \times 10^5 \times 4\times10^{-8}/(0.0165 \times 0.010)
= \qty{33.9}{s^{-2}}$: $T = 2\pi/5.83 = \qty{1.08}{s}$. From $T = \qty{1.10}{s}$:
$\gamma = 4\pi^2mV/(T^2PS^2) = 4\pi^2 \times 0.0165 \times 0.010/(1.21 \times 10^5
\times 4\times10^{-8}) = 1.35$.

\textbf{24.} Friction damps the motion and, if it is dry, shifts the
equilibrium; leakage past the ball lets air escape during a compression,
lowering the restoring force: both lengthen the period and bias
$\gamma$ low --- as the $1.35$ above.

\textbf{25.} Pump and diesel: $Q = 0$, so $\Delta U = W$ --- work became
\omterm{def:b1:kinetic-theory:U}{internal energy} and temperature, by \omterm{prop:b1:first-law:transformations}{Laplace's law} with exponent
$\gamma$. Ball: a fast compression is adiabatic, and $\gamma$ sets the
stiffness of an air spring.
\end{solution}
